% !TEX root = relative_cubic_realization.tex
\documentclass[a4paper,11pt]{ltjsarticle}
\usepackage[margin=25mm,headheight=16pt]{geometry}
\usepackage{amsmath,amssymb,amsthm,mathtools}
\usepackage{enumitem,booktabs}
\usepackage[unicode,hidelinks]{hyperref}
\usepackage{fancyhdr}
\usepackage{microtype}
\setlength{\emergencystretch}{2em}
\allowdisplaybreaks[2]
\pagestyle{fancy}\fancyhf{}
\fancyhead[L]{\small 自己同値付きdg圏のétale局所実現}
\fancyhead[R]{\small 2026-11-13 / v2}
\fancyfoot[C]{\thepage}
\newcommand{\CC}{\mathbb C}
\newcommand{\OO}{\mathcal O}
\newcommand{\Ku}{\operatorname{Ku}}
\newcommand{\Perf}{\operatorname{Perf}}
\newcommand{\RHom}{\operatorname{RHom}}
\newcommand{\otL}{\mathbin{\otimes^{\mathbf L}}}
\newtheorem{theorem}{定理}[section]
\newtheorem{lemma}[theorem]{補題}
\newtheorem{proposition}[theorem]{命題}
\newtheorem{corollary}[theorem]{系}
\theoremstyle{definition}
\newtheorem{definition}[theorem]{定義}
\newtheorem{example}[theorem]{例}
\theoremstyle{remark}
\newtheorem{remark}[theorem]{注意}
\renewcommand{\proofname}{証明}

\hypersetup{pdftitle={自己同値付きdg圏のétale局所実現とArtin切断の限界}}
\title{自己同値付きdg圏のétale局所実現とArtin切断の限界}
\author{}
\date{2026年11月13日／研究草稿 v2}
\begin{document}
\maketitle
\begin{abstract}
自己同値を指定したsmooth properな代数的dg圏の族について、
同値双加群の障害をHochschild cohomologyの不変部分で消去し、
指定された特殊ファイバーの同値をétale局所的に延長する判定法を述べる。
滑らかな三次超曲面のKuznetsov成分への適用では、
幾何的Kodaira--Spencer写像とrotation不変部分の同定を用いる。
また、すべての有限Artin切断が相対smoothであっても、
完備基底上の相対smoothnessは従わないことを、階数四のstrict dg模型で示す。
形式的な障害の消滅、実在する族の相対smoothness、étale局所実現を区別する。
\end{abstract}

\section*{問題設定と結果の射程}
有限型複素スキーム$S$上のsmooth properな代数的dg圏の族とは、
局所的に$\Perf(A)$で表され、$A$が底上perfect、
対角$A$が$A^{\mathrm{op}}\otL A$上perfectとなるものをいう。
自己同値は可逆perfect双加群として与え、圏はpretriangulatedかつ冪等完備とする。
底変更は導来的に取り、同値はdg Morita同値を意味する。
自己同値一つとそのintertwinerを指定するが、三次周期性の選択や
群作用のcoherenceは別のデータとして扱う。

第\ref{er:sec:etale}節では、perfect対象の導来moduli\cite{TV}と
特徴作用による障害\cite{ERLowen}を使い、
自己同値付きの局所実現を定理\ref{er:thm:criterion}で示す。
三次超曲面への適用は定理\ref{er:thm:ER}であり、任意の次元$n\ge3$を対象とする。
ここでは\emph{相対smooth proper性を仮定}する。
次数別flatな量子化からその仮定が自動的に従うとは述べない。
一般点での可換復元定理の証明とは、この入力の水準が異なる。

第\ref{ho:sec:parity}節ではBFKの既知公式\cite{BFKER}から
Hochschildの次数とrotation固有値を整理する。
第\ref{as:sec:obstruction}節では射影直線上の既知の生成性喪失の例\cite{AMSasha}を
strict dg代数として記述し、Artin切断と完備基底の相違を証明する。
この反例のcohomologyは底上flatではないことも明記する。
これらの組合せと一般定式化の先行研究上の初出は未確認である。

\section{同値双加群の障害消去とétale局所実現}\label{er:sec:etale}
本節では、可換模型の存在に使う局所実現原理を証明する。
入力は実在するsmooth properな代数的dg族であり、proper性だけの族ではない。
圏の全変形をHochschild複体から実現するという主張は用いず、
実在するperfect双加群のmoduliと、その完全な障害理論を用いる。
ここでのpairは自己同値一つを指定したもので、$\rho^3\simeq[2]$という
選択された同型やそのcoherenceを追加データとして含めない。

\subsection{perfect加群の小拡大障害}
小拡大は剰余体$k=\CC$を持つ局所Artin代数の全射
$A'\twoheadrightarrow A$で、核$I$が$\mathfrak m_{A'}I=0$を満たすものとする。
$D'$は$A'$上smooth properなdg代数とし、$D=D'\otL_{A'}A$とする。

\begin{lemma}[加群の完全な障害理論と特徴作用]\label{er:lem:obstruction}
$P\in\Perf(D)$の持上げには自然な障害
\[
 o(P;D')\in\operatorname{Ext}^2_{D_k}(P_k,P_k)\otimes_k I
 \tag{O}\label{er:eq:obstruction}
\]
があり、零であることと$\Perf(D')$への持上げの存在は同値である。
障害はMorita関手と対象の同型に関して自然である。
$D'=C'^{\rm op}\otL_{A'}B'$で$P$が$C$から$B$への同値を表すとき、
$B'$を固定し、$C'$のliftをHochschild類$\xi\in\mathrm{HH}^2(C_k)\otimes I$だけ変えると、
障害の変化は、符号を一貫して選べば
\[
 o_{\rm new}-o_{\rm old}=-\chi_P(\xi),\qquad
 \chi_P:\mathrm{HH}^2(C_k)\xrightarrow{\sim}
       \operatorname{Ext}^2_{C_k^{\rm op}\otimes B_k}(P_k,P_k).
 \tag{V}\label{er:eq:variation}
\]
\end{lemma}
\begin{proof}
Toën--Vaquié\cite[Theorem 3.6, Corollaries 3.15, 3.17]{TV}を
可換底$A'$と$D'$に適用する。smoothnessによりpseudo-perfect加群と
perfect加群は一致し、その実際の対象moduliの相対接複体は$\RHom_D(P,P)[1]$である。
導来対象moduliの平方零拡大に関する障害理論は、その接複体の次数一に完全な障害を持つ。
従って障害は$\operatorname{Ext}^2_D(P,P\otL_A I)$に入り、
導来随伴と$\mathfrak m_A I=0$によって\eqref{er:eq:obstruction}に一致する。
これは単なる接空間の計算を代数化の代用にしているのではない。
実在する対象の導来moduliが無限小凝集的であることを使っており、
障害のnull-homotopyはそのmoduliの$A'$-点、すなわち実際のperfect加群を与える。

変化式も説明する。cofibrantな代数模型と加群のbar模型を選び、
構造を次数一のbar余微分$b$で表す。任意のlift$\widetilde b$の曲率
$\widetilde b^{\,2}$は$I$に値を取り、$[b,\widetilde b^{\,2}]=0$である。
持上げを$I$値の$h$だけ変えると、$I^2=0$により曲率は$[b,h]$だけ変わる。
この曲率類が上のmoduliの障害である。
代数側を$\xi$だけ変えた際の曲率の変化は、$\xi$を加群の構造写像に挿入したcochainであり、
それが特徴作用$\chi_P(\xi)$である。符号はsourceの反対代数とtorsorの規約に合わせて
\eqref{er:eq:variation}と取る。必要なのは同じ規約を両式で使うことである。
この記述はLowen\cite[Theorem 4.8]{ERLowen}の特徴作用による障害の記述を、
bar模型でdg加群と任意の小拡大に書いたものである。
$P_k$の逆Morita双加群をテンソルすれば、その自己Homと対角の自己Homが同定されるので、
$\chi_P$は同型である。曲率の構成、またはmoduliの障害の自然性から、最後の自然性も従う。
\end{proof}

\begin{lemma}[還元後の同値とintertwiner]\label{er:lem:lift}
$C',B'$は$A'$上smooth properとする。
$P'\in\Perf(C'^{\rm op}\otL B')$の還元がMorita同値なら$P'$自身も同値である。
さらに$Q'$もperfectで、$\operatorname{Ext}^1(P_k,Q_k)=0$なら、
任意の射$P'_A\to Q'_A$は$A'$上へ持ち上がる。元の射が同型ならliftも同型となる。
\end{lemma}
\begin{proof}
perfect双加群は両側に随伴を持ち、smooth proper性により随伴双加群もperfectである。
unitとcounitのconeは底上perfectで、還元後に零であるから導来Nakayamaで零となる。
射の持上げには三角$Q'\otL I\to Q'\to Q'_A$を$\RHom(P',-)$で評価する。
射の障害は$\operatorname{Ext}^1(P_k,Q_k)\otimes I$に入り、仮定で零である。
持ち上げた射のconeに同じNakayamaを適用すれば同型性も従う。
\end{proof}

\subsection{rotationを保つ障害消去}
$H$は複素数体上smoothな有限型スキーム、$h\in H(\CC)$とする。
$\mathcal A_H$はsmooth proper dg族、$r_H$は可逆perfect双加群で表された自己同値とする。
$\mathcal A=\mathcal A_h$、$r=r_h$と置く。局所的なdg代数模型を取って議論する。

\begin{theorem}[自己同値付き局所実現の判定]\label{er:thm:criterion}
次を仮定する。
\[
 \mathrm{HH}^{<0}(\mathcal A)=0,\qquad
 \mathrm{HH}^{1}(\mathcal A)=0,
 \qquad
 T_hH\longrightarrow\mathrm{HH}^2(\mathcal A)^r
 \text{ が全射}.
 \tag{KS}\label{er:eq:KS}
\]
最後の写像は$\mathcal A_H$のKodaira--Spencer写像である。
$S$は有限型複素スキーム、$(\mathcal B_S,s_S)$はsmooth proper dg族と
可逆perfect双加群のpairとする。$s\in S(\CC)$に指定した同値
$(\mathcal A,r)\simeq(\mathcal B_s,s_s)$があるなら、
pointed étale近傍$V\to(S,s)$、$f:V\to H$、$f(v)=h$と
指定同値を延長するpairの同値
\[
 (f^*\mathcal A_H,f^*r_H)\simeq(\mathcal B_V,s_V)
\]
が存在する。
\end{theorem}
\begin{proof}
まず小拡大$A'\twoheadrightarrow A$上の持上げを証明する。
$\operatorname{Spec}A'\to S$を固定し、$A$上の点$h_A$、同値双加群$P_A$、
intertwiner $\alpha_A:s_AP_A\xrightarrow{\sim}P_Ar_A$が与えられているとする。
$H$はsmoothなので$h_A$の任意のlift$h'_0$を取れる。
両側のliftされた自己同値から、双加群圏の自己同値
\[
 \mathcal F'(P)=s'\circ P\circ(r'_0)^{-1}
\]
を得る。$\alpha_A$は$\mathcal F_A(P_A)\simeq P_A$を与える。
補題\ref{er:lem:obstruction}の自然性より、$P_A$の持上げ障害$o$は
$\mathcal F_k$の作用で不変である。
Morita同定$\operatorname{Ext}^2(P_k,P_k)=\mathrm{HH}^2(\mathcal A)$の下で、
この作用は$r$による共役作用（規約によりその逆）である。
従って
\[
 o\in\mathrm{HH}^2(\mathcal A)^r\otimes I.
 \tag{I}\label{er:eq:invariant}
\]
この段階では$\alpha_A$を$A'$上へ持ち上げたとは仮定していない。

\eqref{er:eq:KS}と\eqref{er:eq:variation}により、$h'_0$を
$T_hH\otimes I$の元だけ修正して$o$を消せる。
修正後のsourceと固定したtargetの間にperfect双加群$P'$が存在し、
補題\ref{er:lem:lift}によって同値となる。
$\alpha_A$を持ち上げる障害は
\[
 \operatorname{Ext}^1(s_kP_k,P_kr_k)\otimes I
       \simeq \mathrm{HH}^1(\mathcal A)\otimes I=0
\]
なので、intertwinerも同型として持ち上がる。
自己同値一つのintertwiner以外に群関係は課していないので、ここで必要なデータは全て揃う。
任意のnilpotent拡大は小拡大に分解できる。

次に代数化する。$H\times S$上で、perfectな
$\mathcal A_H^{\rm op}\otimes\mathcal B_S$-加群$P$と
同型$s_SP\simeq Pr_H$のmoduliを考え、$P$が同値である部分を$\mathcal E$とする。
Toën--Vaquié\cite[Theorem 3.6, Corollaries 3.15, 3.19]{TV}により、
perfect対象とその射の導来moduliは局所幾何的かつ局所有限表示である。
unit、counitおよびintertwinerのconeは底上perfectだから、その消滅は開条件である。
従って$\mathcal E$も局所幾何的かつ局所有限表示である。
同値点での自己Extの負次数は$\mathrm{HH}^{<0}(\mathcal A)=0$であり、
その近傍でこの消滅を保てる。従って古典的切断はその近傍でArtin 1-stackとなる。
ここでは局所dg代数模型を用いてTVの定理を$H\times S$上の可換係数環に適用している。

上の小拡大持上げは、指定点$e=(h,s,P_k,\alpha_k)$で
$\mathcal E\to S$が形式的にsmoothであることを示す。
局所有限表示のArtin stackの無限小判定により、この射は$e$でsmoothである。
$e$を覆うsmooth atlas $U\to\mathcal E$を取り、点の近傍へ縮めると$U\to S$もsmoothである。
滑らかな射のpointed étale局所section\cite[Tags 02HX, 055U]{StacksER}を使い、
$V\to S$と$V\to U\to\mathcal E$を得る。
この普遍対象が所要の族、双加群およびintertwinerを与える。
\end{proof}

\begin{remark}[使わない仮定と先行研究]\label{er:rem:etale-scope}
この証明は$\mathrm{HH}^3=0$、Hochschild cochainのformality、
抽象的なHochschild変形の全てがactual Morita変形になること、または$\rho^3=[2]$の指定を使わない。
Lin--Zhang\cite[Theorem 2.2, Proposition 2.4]{ERLinZhang}の同値核の障害消去と同じ原理を用いるが、
ここではtargetを非幾何的なsmooth proper dg族として扱い、
自己同値が障害を不変部分へ制限することと、intertwinerの持上げを別々に証明した。
同論文の幾何的projectorの仮定を満たさないtargetへその定理を直接適用していない。
この一般定式化の新規性は未確認である。
\end{remark}

\subsection{滑らかなcubicへの適用}
\begin{lemma}[低次Hochschild群と幾何的接写像]\label{er:lem:cubic}
$X=V(F)\subset\mathbb P^{N-1}$を滑らかなcubic、$N=n+2\ge5$とし、
$J(F)=\CC[z_1,\ldots,z_N]/(\partial_1F,\ldots,\partial_NF)$とする。
\[
 \mathrm{HH}^{<0}(\Ku(X))=0,\quad
 \mathrm{HH}^0(\Ku(X))=\CC,\quad \mathrm{HH}^1(\Ku(X))=0,
\]
\[
 \mathrm{HH}^2(\Ku(X))=
 \begin{cases}J(F)_3&n\ne4,\\J(F)_3\oplus\CC^2&n=4,\end{cases}
 \qquad
 H^1(X,T_X)\xrightarrow{\sim}\mathrm{HH}^2(\Ku(X))^{\rho_X}=J(F)_3.
 \tag{H2}\label{er:eq:HH}
\]
$\dim J(F)_3=\binom N3$である。
\end{lemma}
\begin{proof}
BFK\cite[Theorem 5.39]{BFKER}のequivariant matrix factorizationの公式を用いる。
$G=\mathbb G_m$、各座標の重みは1、potentialのcharacterは3である。
恒等sectorのKoszul cohomologyは、偏微分が正則列であるためJacobian商だけで、
cohomological degree $2r$に$J(F)_{3r}$を与える。
非恒等$g\in\mu_3$の固定空間は原点、conormalの次元は$N$、determinantの重みは$-N$である。
次数$q$の重み零条件は$(3q-N)/2=0$であり、従って$3\mid N$かつ$q=N/3$のときだけ
各sectorから一次元が現れる。これで負次数、次数0、1、2の群が得られる。

幾何的Kodaira--Spencer写像の同定は、単なる次元一致ではない。
$n\ge4$では$3\le(n+2)/2$なので、Huybrechts--Rennemo
\cite[Lemma 2.19, \S2.5, \S3.3]{HRER}のAtiyah--Kodaira--Spencerと
三次周期性の比較を適用でき、$F+\epsilon G$の像は恒等sectorの$[G]\in J(F)_3$である。
$n=3$ではJacovskis--Lin--Liu--Zhang\cite[Theorem 4.4(2)]{JLLZER}の
幾何的写像の同型を用いる。normal sequenceとEuler sequenceから$H^1(T_X)=J(F)_3$である。
幾何的族にはrotationがあるのでこの像は不変である。
$n=4$で余分な二方向が不変でないことは、Huybrechts\cite[Proposition 3.12]{HuyER}による。
rotationはsymplecticで、Mukai格子の$A_2$上で三つの根を巡回する。
不変なCalabi--Yau類による$\mathrm{HH}^2\simeq\mathrm{HH}_0$の同定の下で、
その二方向の固有値は$\zeta,\zeta^2$である。よって不変部分は正確に$J(F)_3$となる。
最後に偏微分が二次正則列であることから$H_{J(F)}(t)=(1+t)^N$を得る。
\end{proof}

\begin{theorem}[cubicのrotation付きétale局所実現]\label{er:thm:ER}
$S$は有限型複素スキーム、$\mathcal K_S$は実在するsmooth properな代数的dg族で、
$\rho_S$は可逆perfect双加群により表される自己同値とする。
$s\in S(\CC)$で、$n\ge3$の滑らかなcubic $X$に対する指定同値
$(\mathcal K_s,\rho_s)\simeq(\Ku(X),\rho_X)$があるなら、
pointed étale近傍$V\to(S,s)$、滑らかなcubic族$Y/V$、および指定同値を延長する同値
\[
 (\mathcal K_V,\rho_V)\simeq(\Ku(Y/V),\rho_{Y/V})
 \tag{ER}\label{er:eq:ER}
\]
が存在する。
\end{theorem}
\begin{proof}
滑らかな三次形式の射影パラメータ空間の$[F]$を含むアフィン開を$H$とする。
普遍cubic $\mathcal X/H$のrelative perfect圏はsmooth properで、
$\OO,\ldots,\OO(n-2)$のrelative exceptional列の射影により
smooth properな相対Kuznetsov成分を得る。
射影とrotationは有限個の評価三角で構成され、底変更に可換である。
普遍cubic上の有限個のtwistから生成子を選べるので、局所的なdg代数模型も存在する。
$T_{[F]}H\to H^1(T_X)=J(F)_3$は全射である。
補題\ref{er:lem:cubic}と定理\ref{er:thm:criterion}を適用する。
\end{proof}

\begin{corollary}[rotationを最初に指定しない存在]\label{er:cor:unpolarized}
$n\ne4$なら、前定理で$\rho_S$を入力として指定せずとも、
$\mathcal K_s\simeq\Ku(X)$からétale局所的な可換cubic模型が得られる。
これはrotationなしの大域的な一意性の主張ではない。
\end{corollary}
\begin{proof}
補題\ref{er:lem:cubic}では幾何的接写像が$\mathrm{HH}^2$全体に全射である。
定理\ref{er:thm:criterion}の証明から不変性とintertwinerの段階を省けばよい。
\end{proof}

\section{Hochschildの奇数次障害とrotation不変部分}\label{ho:sec:parity}
可換復元における対角smoothnessを、形式変形の障害と混同しないために、
滑らかな三次超曲面のKuznetsov成分の低次Hochschild群を全次元で整理する。
以下は既知のequivariant matrix factorizationの計算からの帰結であり、
残余dg圏の相対smoothnessに関する新たな一般定理ではない。

\begin{proposition}[奇数次Hochschild群とrotation固有値]\label{ho:prop:sectors}
$X=V(F)\subset\mathbb P^{N-1}_\CC$を滑らかな三次超曲面、$N=n+2\ge5$とし、
$J(F)=\CC[z_1,\dots,z_N]/(\partial F)$とする。
$\operatorname{Ku}(X)$の標準rotationを$\rho$とする。
$\zeta$を原始三乗根と書けば、次の複素ベクトル空間の分解が成り立つ：
\[
 \operatorname{HH}^{q}(\operatorname{Ku}(X))\cong
 \begin{cases}
 J(F)_{3q/2},&q\text{ 偶数},\\ 0,&q\text{ 奇数}
 \end{cases}
 \quad\oplus\quad
 \begin{cases}
 \CC_{\zeta}\oplus\CC_{\zeta^2},&3\mid N,\ q=N/3,\\
 0,&\text{それ以外}.
 \end{cases}
 \tag{HS}\label{ho:eq:sectors}
\]
ここで$J(F)_d=0$（$d<0$）とし、右側の二つの一次元空間は
$\rho$の共役作用の非自明な固有空間である。
$J(F)$から来る項では$\rho$は恒等的に作用する。
特に
\[
 \operatorname{HH}^3(\operatorname{Ku}(X))
 \cong\begin{cases}\CC^2&n=7,\\0&n\ne7,\end{cases}
 \qquad
 \boxed{\operatorname{HH}^3(\operatorname{Ku}(X))^\rho=0\quad(n\ge3).}
 \tag{V3}\label{ho:eq:HH3}
\]
さらに$\operatorname{HH}^2(\operatorname{Ku}(X))^\rho\cong J(F)_3$で、
$\dim_\CC J(F)_3=\binom N3$である。
\end{proposition}
\begin{proof}
Orlovのgraded matrix factorization表示と
Ballard--Favero--Katzarkov\cite[Theorem 5.39, pp.~82--85]{BFKER}の
固定sector公式を用いる。
$G=\mathbb G_m$は全生成変数の重みが$1$となる作用で、
potentialのcharacterは三次である。$K_\chi=\mu_3$とする。
$g=1$において$V_g=\mathbb A^N$、$W_g=0$である。
偏微分$(\partial F)$は原点にのみ共通零点を持つ長さ$N$の正則列であり、
Koszul cohomologyは零次の$J(F)$に集中する。
重み零の条件とcohomological gradingから、恒等sectorは偶数次数$q=2r$に
$J(F)_{3r}$を与え、奇数次数には寄与しない。

$g\in\mu_3\setminus\{1\}$なら固定空間は原点で、$\dim W_g=N$、
$J(F_g)=\CC$である。
Theorem 5.39のtwistとdeterminant characterを代入すると、
$u=(q-N)/2$に対する重み零条件は$3u+N=0$となる。
従って$q=N/3$でのみ各非自明sectorが一次元を与える。
各sectorは標準的な文字twist $(1)$ による共役作用で
$g$の文字（規約によって逆文字）を受ける。
これは同論文のkernelの固定graph分解と
\cite[Theorem 6.5]{BFKER}のcharacter作用の計算と同じものである。
Orlov同値の下で次数shift $(1)$ は$\rho$に対応するため、
二つのsectorの固有値は順序を除いて$\zeta,\zeta^2$となる。
$J(F)$から来る恒等sectorではこの作用は自明である。
これで\eqref{ho:eq:sectors}を得る。

$q=3$は奇数なので、非零になり得るのは$N/3=3$、すなわち$N=9$
（$n=7$）だけである。しかもその二方向には$\rho$-不変ベクトルがない。
$q=2$の不変部分は常に$J(F)_3$である。
滑らかさから偏微分は二次正則列なので、
$H_{J(F)}(t)=(1+t)^N$となり$\dim J(F)_3=\binom N3$を得る。
\end{proof}

\begin{corollary}[非楕円cubic fivefoldの変形障害の所在]\label{ho:cor:fivefold}
滑らかな三次五次元超曲面$X^5$について
（特に節点型・尖点型の量子化で現れる滑らかな可換特殊ファイバーにも）、
\[
 (\dim\operatorname{HH}^0,\dim\operatorname{HH}^1,
 \dim\operatorname{HH}^2,\dim\operatorname{HH}^3,
 \dim\operatorname{HH}^4)=(1,0,35,0,7).
 \tag{F5}\label{ho:eq:fivefold}
\]
従って滑らかな可換特殊ファイバーのdg圏自身の通常のMorita変形に対する
\emph{一次の}Kuranishi障害空間$\operatorname{HH}^3$は零である。
ただし、量子化から構成した$\Gamma_R$の
$\Gamma_R^e$上の対角perfect性を、これだけからは結論できない。
\end{corollary}
\begin{proof}
$N=7$を命題\ref{ho:prop:sectors}に代入する。
$J(F)_0,J(F)_3,J(F)_6$の次元はそれぞれ$1,35,7$で、
それ以外の表示次数$q=1,3$は零となる。
Morita変形の一次障害が$\operatorname{HH}^3$に入ることは
標準的なHochschild変形理論の適用である。
\end{proof}

\begin{remark}[残余対角との論理的境界]\label{ho:rem:barrier}
$\operatorname{HH}^3(\Ku(X))^\rho=0$は重要な形式的情報であるが、
この不変部分の消滅だけで任意のrotation付きdg圏変形の厳密な
$\mathbb Z/3$-equivariant obstruction theoryを構成したとはしない。
また、Artin切断上の変形が可換的に説明できたとしても、
有限型係数環上の\emph{実在する}残余対角のperfect性、
すなわち$\Gamma_R\in\Perf(\Gamma_R^{\mathrm{op}}\otL_R\Gamma_R)$は別の問題である。
具体的な量子化族へ定理\ref{er:thm:ER}を適用する際には、
この残余対角のperfect性を別途確認しなければならない。
本節はBFKの既知の公式の具体的な帰結と、その本研究における障害の切分けを与える。
\end{remark}

% 2026-11-10 / v22
\section{全てのArtin切断のsmoothnessから完備基底のsmoothnessは従わない}
\label{as:sec:obstruction}

本節では、有限Artin基底での生成子の持上げと、
完備局所基底での生成性の間にある障害を明示する。
出発点となる$\mathbb P^1$上の生成対象の変形例は、
SashaのMathOverflow回答\cite{AMSasha}にある既知の例である。
以下の四次元strict dg模型、および全Artin切断の相対smoothnessと
generic fiberの非smoothnessを併記した整理の先行研究上の初出は未確認である。

この例は、非可換cubicから構成した本稿の$\Gamma_R$が非smoothであることを示すものではない。
示すのは、底上proper、滑らかな特殊ファイバー、
全ての有限Artin切断の相対smoothnessという抽象的条件だけでは、
完備基底上の相対smoothnessを結論できないという事実である。

\subsection{生成対象がgeneric fiberで生成しなくなる族}

$R=\CC[[h]]$、$K=\CC((h))$、$X=\mathbb P^1_R$とし、
$H^0(X,\mathcal O_X(1))$の基底を$s,t$とする。
ここで$s,t$は射影直線の一次sectionであり、
MFの周期変数$t$とは無関係である。
混同を避けるため、本節でのみこの記号を用いる。
二項perfect複体
\begin{equation}\label{as:eq:P}
 P_h=\bigl[\mathcal O_X(-1)\xrightarrow{\,hs\,}\mathcal O_X\bigr]
\end{equation}
をcohomological degrees $-1,0$に置く。

特殊ファイバーでは微分が零なので
\[
 P_0\simeq\mathcal O_{\mathbb P^1}\oplus
                  \mathcal O_{\mathbb P^1}(-1)[1].
\]
これは$\operatorname{Perf}(\mathbb P^1)$の古典的生成対象である。
一方、generic fiberでは$h$は可逆なので、
$s$の零点$p$について
\[
 P_h\otimes_R K\simeq\mathcal O_p.
\]
$\mathcal O_p$から有限coneと直和因子によって作る対象のsupportは$p$に含まれるため、
$\mathcal O_p$は$\operatorname{Perf}(\mathbb P^1_K)$を生成しない。
従って$P_h$も$\operatorname{Perf}(\mathbb P^1_R)$を生成しない。
実際、生成していれば、その有限生成操作を$K$へ底変更することで
generic fiberでも生成してしまう。

\begin{lemma}[各有限切断では生成する]\label{as:lem:artin-generation}
$R_N=R/(h^{N+1})$とし、$P_N=P_h\otimes_RR_N$とする。
このとき$P_N$は$\operatorname{Perf}(\mathbb P^1_{R_N})$を生成する。
\end{lemma}
\begin{proof}
射影直線の相対Beilinson分解により、perfect圏の底変更と
perfect対象間のHomの導来底変更は$R_N\to\CC$と両立する。
従って補題\ref{am:lem:generator}のInd圏における議論を、そのまま
$\operatorname{Perf}(\mathbb P^1_{R_N})$に適用できる。
具体的には、$P_N$と右直交するInd対象$M$について、
$P_N$のcompact性と底変更から
\[
 \operatorname{RHom}(P_0,M\otimes_{R_N}^{\mathbf L}\CC)=0
\]
となる。$P_0$は生成対象なので$M\otimes_{R_N}^{\mathbf L}\CC=0$である。
$R_N$の有限$h$-filtrationにより
$R_N\in\operatorname{thick}_{D(R_N)}(\CC)$だから、$M=0$となる。
これで$P_N$のcompact生成性、従ってperfect圏の古典的生成性を得る。

ここで用いた射影直線の有限模型は、対角の完全列
\[
 0\longrightarrow\mathcal O(-1,-1)
 \longrightarrow\mathcal O
 \longrightarrow\mathcal O_\Delta\longrightarrow0
\]
から得られる。従ってこの底変更の主張は、非正則なArtin基底でも成立する。
\end{proof}

\subsection{四次元strict dg代数}

\begin{proposition}[底上properな明示的dg模型]\label{as:prop:dg-model}
$\Gamma_R=\operatorname{RHom}_X(P_h,P_h)$は、次のstrict unital dg代数で表される：
\begin{equation}\label{as:eq:dg-algebra}
 \Gamma_R^0=Re_-\oplus Re_+,\qquad
 \Gamma_R^1=Rs\oplus Rt,\qquad
 \Gamma_R^q=0\quad(q\ne0,1),
\end{equation}
\[
 d(e_-)=hs,\qquad d(e_+)=-hs,\qquad d(s)=d(t)=0.
\]
$e_-,e_+$は互いに直交する冪等元で、単位は$e_-+e_+$である。
$s,t$は$e_-$から$e_+$へのdegree一の射であり、
\[
 e_+s=s=se_-,\qquad e_+t=t=te_-,
\]
その他の冪等元と射の積、および二本の射同士の積は零である。
特に$\Gamma_R$は$R$上perfectである。
この模型の形成は各$R_N$と$K$への導来底変更と両立する。
\end{proposition}
\begin{proof}
$P_h$の二つの項の間のsheaf Homは
$\mathcal O_X(-1),\mathcal O_X,\mathcal O_X(1)$の直和である。
これらの一次以上のcohomologyは零なので、
大域Homの二項複体が$\operatorname{RHom}$を実際に計算する。
degree $-1$の大域Homは$H^0(\mathcal O_X(-1))=0$であり、
degree零は二つの項のscalar endomorphism、degree一は
$H^0(\mathcal O_X(1))=Rs\oplus Rt$である。
複体の微分を$D=hs$とすると、endomorphismの微分は
$f\mapsto Df-(-1)^{|f|}fD$なので、表示した微分を得る。
積はendomorphismの合成そのものである。

この計算は有限自由$R$-加群とその積だけからなり、
同じline bundle cohomologyの消滅は各底変更後も成立する。
従って通常のtensorがこの模型の導来底変更を計算し、
底変更後の$P_h$のderived endomorphismを与える。
\end{proof}

\begin{lemma}[generic fiberの外積代数模型]\label{as:lem:generic}
$E=K[\epsilon]/(\epsilon^2)$に$|\epsilon|=1$、零微分を与える。
このとき
\[
 E\longrightarrow\Gamma_R\otimes_R K,\qquad
 1\longmapsto e_-+e_+,\qquad \epsilon\longmapsto t
\]
はdg代数のquasi-isomorphismである。
\end{lemma}
\begin{proof}
単位と積と微分を保つことは式\eqref{as:eq:dg-algebra}から直ちに分かる。
$h$が可逆なので
\[
 H^0(\Gamma_R\otimes_R K)=K(e_-+e_+),\qquad
 H^1(\Gamma_R\otimes_R K)=(Ks\oplus Kt)/Ks=K\cdot[t]
\]
であり、他のcohomologyは零である。
従って上の写像は全てのcohomologyで同型となる。
\end{proof}

\subsection{無限derived tensorによる非smoothness}

\begin{lemma}[degree一の外積代数はsmoothでない]
\label{as:lem:exterior-not-smooth}
前補題の$E$は$K$上properだが、homologically smoothではない。
\end{lemma}
\begin{proof}
augmentation $E\to K$によって$K$を左右の$E$-加群とする。
左$E$-加群のsemifree resolution
\[
 Q=\bigoplus_{n\ge0}Eu_n,\qquad |u_n|=0,\qquad
 d(u_0)=0,\quad d(u_n)=\epsilon u_{n-1}\quad(n\ge1)
\]
を考える。$Q\to K$は$u_0\mapsto1$、$u_n\mapsto0$（$n\ge1$）で与える。
$n$に関する増大filtrationはsemifree filtrationであり、
$d^2=0$は$\epsilon^2=0$による。
degree零のcycleは$Ku_0$だけであり、degree一の全ての元
$\epsilon u_n$は$d(u_{n+1})$である。
従って$Q\to K$はquasi-isomorphismである。

右からaugmentation加群$K$をtensorすると
\[
 K\otimes_E^{\mathbf L}K
 \simeq K\otimes_E Q
 =\bigoplus_{n\ge0}Ku_n
\]
となり、右辺の微分は零で、全ての$u_n$はdegree零にある。
従って$H^0(K\otimes_E^{\mathbf L}K)$は無限次元である。
もし$K$が$E$上perfectなら、このderived tensorは
有限個の$K$のshiftから有限coneと直和因子で作られ、
有限次元の有界cohomologyを持つはずなので矛盾する。

さらに、もし$E$の対角が$E^{\rm op}\otimes_KE$上perfectなら、
その有限cell表示の一方の$E$-作用へ$-\otimes_E^{\mathbf L}K$を施すと、
$K$の有限$E$-cell表示を得る。
これは$K$が$K$上perfectであることを用いる。
従ってsmoothnessは$K\in\operatorname{Perf}(E)$を含意するが、
直前でこれを否定した。
よって$E$はhomologically smoothでない。
一方、$E$自体は二次元の有界$K$-複体なのでproperである。
\end{proof}

\begin{theorem}[全Artin切断ではsmoothだが完備基底ではsmoothでない族]
\label{as:thm:obstruction}
命題\ref{as:prop:dg-model}の$\Gamma_R$は、次の全ての性質を持つ。
\begin{enumerate}
 \item $\Gamma_R$は$R$上perfectなstrict dg代数である。
 \item $\Gamma_R\otimes_R^{\mathbf L}\CC$は$\CC$上smoothである。
 \item 各$N\ge0$に対し
 $\Gamma_R\otimes_R^{\mathbf L}R_N$は$R_N$上smoothである。
 \item $\Gamma_R\otimes_R K$は$K$上smoothでなく、
 $\Gamma_R$も$R$上smoothでない。
\end{enumerate}
\end{theorem}
\begin{proof}
第一項は有限自由模型から従う。
補題\ref{as:lem:artin-generation}により$P_N$は
$\operatorname{Perf}(\mathbb P^1_{R_N})$の生成対象であり、
命題\ref{as:prop:dg-model}の底変更から
\[
 \operatorname{Perf}(\Gamma_R\otimes_R^{\mathbf L}R_N)
 \simeq\operatorname{Perf}(\mathbb P^1_{R_N})
\]
となる。
相対射影直線の対角は前述の二項線束分解を持つので、
このperfect圏は$R_N$上smoothである。
smoothnessはdg Morita同値で不変だから、第二・第三項を得る。

補題\ref{as:lem:generic}と
\ref{as:lem:exterior-not-smooth}からgeneric fiberはsmoothでない。
$\Gamma_R$が$R$上smoothなら、そのperfect対角の表示を
flatな底変更$R\to K$へ移してgeneric fiberもsmoothとなるため、矛盾する。
従って第四項も成立する。
\end{proof}

\begin{remark}[この反例の適用範囲]
$R_N\in\operatorname{thick}_{D(R_N)}(\CC)$はArtin基底での生成子の持上げを支える。
しかし$R=\CC[[h]]$は$D(R)$で$\CC$の有限coneと直和因子から作れない。
実際、そのような対象を$K$へtensorすると零になるが、$R\otimes_RK=K\ne0$である。
上の族では、この違いが生成性とsmoothnessの実際の破綻として現れる。

従って、全Artin比較の両立性だけから、完備基底上の生成性や相対smoothnessを補ってはならない。
可換側がsmooth properな射影族であることだけでも、変形した対象の生成性は保証されない。
中心三次非可換商の残余圏へこの原理を適用するには、
当該windowと生成対象についての有限生成証明、
または対角の具体的なperfect表示を別途用いる。
本例は、その特別な主張を否定してはいない。

なお本例の$R$上のcohomologyは
$H^0(\Gamma_R)=R$、$H^1(\Gamma_R)=Rt\oplus(R/hR)s$であり、
cohomologyの底上flatnessは満たさない。
この例が否定する仮定の範囲を、証明していない別の仮定へ広げてはならない。
例\ref{cg:ex:artin-singularity}はMFと絶対特異圏の区別を示すものであり、
本節のsmoothnessの持上げ障害とは別の反例である。
\end{remark}

\begin{thebibliography}{99}
\bibitem{TV} B. Toën and M. Vaquié,
\emph{Moduli of objects in dg-categories},
Ann. Sci. École Norm. Sup. (4) \textbf{40} (2007), 387--444.
\url{https://doi.org/10.1016/j.ansens.2007.05.001}.
Theorem 3.6, Corollaries 3.15, 3.17, 3.19, 3.21。
\url{https://www.math.univ-toulouse.fr/~vaquie/marque/vaquie_AnnEns40.pdf}.

\bibitem{ERLowen} W. Lowen,
\emph{Hochschild cohomology, the characteristic morphism and derived deformations},
Compos. Math. \textbf{144} (2008), 1557--1580.
\url{https://arxiv.org/abs/0707.2602v1}.
Theorem 4.8, Section 4.4。任意の小拡大に対するdg加群の記述は本文で説明する。

\bibitem{BFKER} M. Ballard, D. Favero and L. Katzarkov,
\emph{A category of kernels for equivariant factorizations and its implications for Hodge theory},
Publ. Math. Inst. Hautes Études Sci. \textbf{120} (2014), 1--111.
\url{https://doi.org/10.1007/s10240-013-0059-9}.
Theorems 5.39, 6.5。

\bibitem{HRER} D. Huybrechts and J. V. Rennemo,
\emph{Hochschild cohomology versus the Jacobian ring and the Torelli theorem for cubic fourfolds},
Algebraic Geometry \textbf{6} (2019), 76--99.
\url{https://doi.org/10.14231/AG-2019-005}.
Lemma 2.19, Sections 2.5, 3.3。

\bibitem{JLLZER} A. Jacovskis, X. Lin, Z. Liu and S. Zhang,
\emph{Infinitesimal categorical Torelli theorems for Fano threefolds},
J. Pure Appl. Algebra \textbf{227} (2023), 107418.
\url{https://arxiv.org/abs/2203.08187}.
Theorem 4.4(2)。

\bibitem{HuyER} D. Huybrechts,
\emph{The K3 category of a cubic fourfold},
Compos. Math. \textbf{153} (2017), 586--620.
\url{https://doi.org/10.1112/S0010437X16008137}.
Proposition 3.12, Section 3.5。

\bibitem{StacksER} The Stacks Project Authors, \emph{The Stacks Project}.
\url{https://stacks.math.columbia.edu/tag/02HX};
\url{https://stacks.math.columbia.edu/tag/055U}.
2026年11月13日確認。smoothnessの無限小判定とsmooth射のétale局所section。

\bibitem{ERLinZhang} X. Lin and S. Zhang,
\emph{Categorical Torelli and Fourier--Mukai partners of hypersurfaces},
arXiv:2611.05915v1 (2026).
\url{https://arxiv.org/abs/2611.05915}.
Section 2.1, Theorem 2.2, Proposition 2.4（元原稿の参照箇所）。
本稿の判定法の証明の入力ではなく、関連研究として挙げる。

\bibitem{AMSasha} Sasha,
Answer to \emph{Is the property of being a dg generator open?},
MathOverflow, 5 December 2016.
\url{https://mathoverflow.net/questions/256391/is-the-property-of-being-a-dg-generator-open}.
射影直線上の生成対象がskyscraperへ変形して生成性を失う例。

\end{thebibliography}
\end{document}
